%% 
%% Copyright 2019-2024 Elsevier Ltd
%% 
%% This file is part of the 'CAS Bundle'.
%% --------------------------------------
%% 
%% It may be distributed under the conditions of the LaTeX Project Public
%% License, either version 1.3c of this license or (at your option) any
%% later version.  The latest version of this license is in
%%    http://www.latex-project.org/lppl.txt
%% and version 1.3c or later is part of all distributions of LaTeX
%% version 1999/12/01 or later.
%% 
%% The list of all files belonging to the 'CAS Bundle' is
%% given in the file `manifest.txt'.
%% 
%% Template article for cas-dc documentclass for 
%% double column output.

\documentclass[a4paper,fleqn]{cas-dc}

% If the frontmatter runs over more than one page
% use the longmktitle option.

%\documentclass[a4paper,fleqn,longmktitle]{cas-dc}

%\usepackage[numbers]{natbib}
%\usepackage[authoryear]{natbib}
\usepackage[authoryear,longnamesfirst]{natbib}

%%%Author macros
\def\tsc#1{\csdef{#1}{\textsc{\lowercase{#1}}\xspace}}
\tsc{WGM}
\tsc{QE}
%%%

% Uncomment and use as if needed
%\newtheorem{theorem}{Theorem}
%\newtheorem{lemma}[theorem]{Lemma}
%\newdefinition{rmk}{Remark}
%\newproof{pf}{Proof}
%\newproof{pot}{Proof of Theorem \ref{thm}}

\begin{document}
\let\WriteBookmarks\relax
\def\floatpagepagefraction{1}
\def\textpagefraction{.001}

% Short title
\shorttitle{}    

% Short author
\shortauthors{}  

% Main title of the paper
\title []{Effectiveness of Cross Project Learning in Bug Localization}  

% Title footnote mark
% eg: \tnotemark[1]
\tnotemark[1] 

% Title footnote 1.
% eg: \tnotetext[1]{Title footnote text}
\tnotetext[1]{} 

% First author
%
% Options: Use if required
% eg: \author[1,3]{Author Name}[type=editor,
%       style=chinese,
%       auid=000,
%       bioid=1,
%       prefix=Sir,
%       orcid=0000-0000-0000-0000,
%       facebook=<facebook id>,
%       twitter=<twitter id>,
%       linkedin=<linkedin id>,
%       gplus=<gplus id>]

\author[1]{A Eashaan Rao}%[<options>]


% Corresponding author indication
% \cormark[1]

% Footnote of the first author
% \fnmark[1]

% Email id of the first author
\ead{cs21d002@iittp.ac.in}

% URL of the first author
% \ead[url]{}

% Credit authorship
% eg: \credit{Conceptualization of this study, Methodology, Software}
% \credit{}

% Address/affiliation
\affiliation[1]{organization={Indian Institute of Technology Tirupati},
            addressline={}, 
            city={Tirupati},
%          citysep={}, % Uncomment if no comma needed between city and postcode
            postcode={517619}, 
            state={Andhra Pradesh},
            country={India}}

\author[2]{Sridhar Chimalakonda}%[]

% Footnote of the second author
% \fnmark[2]

% Email id of the second author
\ead{ch@iittp.ac.in}

% URL of the second author
% \ead[url]{}

% Credit authorship
% \credit{}

% Address/affiliation
\affiliation[2]{organization={Indian Institute of Technology Tirupati},
            % addressline={}, 
            city={Tirupati},
%          citysep={}, % Uncomment if no comma needed between city and postcode
            postcode={517619}, 
            state={Andhra Pradesh},
            country={India}}

% Corresponding author text
\cortext[1]{Corresponding author}

% Footnote text
% \fntext[1]{}

% For a title note without a number/mark
%\nonumnote{}

% Here goes the abstract
\begin{abstract}
The literature currently lacks a comprehensive empirical analysis that investigates the factors governing the success or failure of cross-project learning for bug localization. While the potential of CPL is acknowledged, there is little to no guidance on its practical application. The investigation is designed to move beyond a binary assessment of CPL's efficacy. Instead, it seeks to answer the critical questions of when, why, and how CPL is effective. This study systematically evaluate the performance of CPL for SOTA's BL models across a diverse set of software projects. In addition, identify and analyze the key factors (e.g., project similarity, dataset size, domain characteristics) that correlate with successful or failed knowledge transfer with support of quantify the impact of negative transfer and identify the scenarios where it is most likely to occur.

The primary outcome of this study will be a set of empirically-grounded, actionable guidelines for researchers and practitioners. These guidelines will detail how to select appropriate source projects and configure training settings to maximize the effectiveness and avoid the pitfalls of future CPL-based BL models.
\end{abstract}

% Use if graphical abstract is present
%\begin{graphicalabstract}
%\includegraphics{}
%\end{graphicalabstract}

% Research highlights
% \begin{highlights}
% \item 
% \item 
% \item 
% \end{highlights}


%\nocite{*}

% Keywords
% Each keyword is seperated by \sep
\begin{keywords}
 \sep \sep \sep
\end{keywords}

\maketitle

% Main text
\section{Introduction}
he primary objective of this phase of research is to conduct a systematic and comprehensive empirical analysis of Cross-Project Learning (CPL) for the task of bug localization. The core difficulty lies in bridging the semantic gap between the natural language of bug reports and the structured nature of source code. While many methods have been developed, their performance remains suboptimal for practical use, a problem exacerbated by the scarcity of large, high-quality datasets for training robust models. CPL has emerged as a promising technique to address this by leveraging data from multiple source projects to improve performance on a target project. However, CPL introduces its own challenges, such as domain shift and vocabulary mismatch, and its effectiveness has not yet been thoroughly studied under controlled, multi-metric evaluation conditions. This study addresses that gap by not only evaluating whether CPL works, but by systematically identifying the conditions under which it succeeds or fails, and by empirically validating the first data-driven strategy for source project selection.

To achieve this, the study employs \textbf{TRANP-CNN} \cite{huo2019deep}, a convolutional neural network-based reranking model, as the representative SOTA bug localization model for Phase 1 evaluation. The evaluation is structured around four training scenarios designed to isolate the contribution of cross-project data at varying levels of target-project data availability. The study is guided by the following research questions:

\begin{itemize}
    \item \textbf{RQ1}: Does cross-project transfer (CP-transfer) significantly outperform within-project training on limited data (WP-small) across multiple evaluation metrics?\\
    \textbf{Motivation}: This is the core CPL feasibility question. Prior work has asserted CPL's promise but has not reported matched-budget comparisons with statistical significance testing across a large and diverse project corpus.

    \item \textbf{RQ2}: What target-project characteristics are the strongest predictors of CPL performance, and do these predictors hold consistently across retrieval and ranking metrics?\\
    \textbf{Motivation}: Understanding which project properties drive or limit performance is essential for practical deployment guidance. This question investigates whether source- or target-side features are more predictive, and whether recall and ranking metrics reveal different signals.

    \item \textbf{RQ3}: Can a data-driven heuristic reliably identify the best source project for a given target, and what is the upper bound on source selection quality?\\
    \textbf{Motivation}: No prior CPBL study has formally validated source selection strategies beyond intuitive domain-matching. This question establishes an empirical benchmark for source selection quality using Kendall rank correlation and Hit@1 accuracy.

    \item \textbf{RQ4}: Is the FAISS-based retrieval stage a limiting factor in the end-to-end CPL pipeline, or does the performance ceiling reside entirely in the reranking model?\\
    \textbf{Motivation}: Diagnosing whether failures originate in retrieval (ground-truth file not retrieved) or in reranking (file retrieved but ranked too low) directly informs where architectural investment is most valuable.
\end{itemize}

\section{Experimental Setup}
\paragraph{Dataset Selection and Curation.}
The candidate pool for this study was assembled by consolidating projects from five established bug localization benchmark datasets: \textbf{Bench4BL} \cite{lee2018bench4bl}, \textbf{BeetleBox} \cite{chakraborty2025blaze}, \textbf{LCA} \cite{bogomolov2024long}, \textbf{SWE-bench} \cite{jimenez2023swe}, and the dataset of \textbf{Ye et al.} \cite{ye2014learning}. Across these five sources, the initial candidate pool comprised \textbf{over 500 unique repositories} spanning multiple programming languages. To ensure that experiments are conducted on projects with statistically adequate bug report volumes, projects were filtered according to language-stratified thresholds: for the primary languages (Python and Java), a minimum of \textbf{100 unique bug reports} per project was required; for under-represented languages (Go, C++, JavaScript, Kotlin), a relaxed threshold of \textbf{at least 20 bug reports} was applied to retain sufficient cross-language diversity. After deduplication and quality filtering, this process yielded a curated multi-language pool of \textbf{98 repositories} covering 41 Python, 34 Java, 9 Kotlin, 6 JavaScript, 5 C++, and 3 Go projects.

For Phase 1 of the Objective 1 evaluation, the scope is restricted to \textbf{14 Python projects} to establish a clean, language-homogeneous baseline for CPL analysis. This restriction avoids confounding the CPL signal with cross-language transfer effects, which will be investigated separately in Objective 2. The 14 projects were selected to ensure diversity in codebase size (39K--696K LoC), application domain (scientific computing, machine learning, data visualization, developer tools, cloud infrastructure), and historical bug report volume (100--789 reports). This produces \textbf{87 unique source$\rightarrow$target project pairs} (all ordered pairs over 14 projects minus same-project pairs where data permits), yielding \textbf{345 individual experiment runs} across four evaluation scenarios. Table~\ref{tab:dataset} summarizes the 14 projects.

\begin{table*}[h]
\centering
\caption{The 14 Python projects used in Phase 1, ordered by codebase size.}
\label{tab:dataset}
\begin{tabular}{|l|c|c|c|c|}
\hline
\textbf{Project} & \textbf{LoC} & \textbf{Bug Reports} & \textbf{Domain} & \textbf{Size Category} \\ \hline
jupyterlab/jupyterlab   &  39,372 & 197 & Developer Tools      & Small  \\ \hline
google/jax              &  72,684 & 326 & Scientific/ML        & Small  \\ \hline
lightning-ai/lightning  &  51,387 & 366 & Scientific/ML        & Small  \\ \hline
prefecthq/prefect       & 107,073 & 404 & Cloud Infrastructure & Medium \\ \hline
pydata/xarray           & 142,044 & 110 & Scientific Computing & Medium \\ \hline
rucio/rucio             & 156,628 & 404 & Cloud Infrastructure & Medium \\ \hline
open-mmlab/mmdetection  & 208,928 & 102 & Computer Vision      & Large  \\ \hline
matplotlib/matplotlib   & 248,837 & 183 & Data Visualization   & Large  \\ \hline
ray-project/ray         & 244,283 & 337 & Scientific/ML        & Large  \\ \hline
numpy/numpy             & 276,548 & 789 & Scientific Computing & Large  \\ \hline
scikit-learn/scikit-learn & 376,169 & 228 & Machine Learning    & XLarge \\ \hline
scipy/scipy             & 438,217 & 100 & Scientific Computing & XLarge \\ \hline
sympy/sympy             & 696,129 & 384 & Mathematics          & XLarge \\ \hline
dmwm/wmcore             & 1,048,906 & 277 & Cloud Infrastructure & XLarge \\ \hline
\end{tabular}
\end{table*}

\subsection{Dataset Preprocessing}
The preprocessing pipeline is structured into two distinct stages. In the first stage, bug report preprocessing, raw bug reports are extracted from the respective benchmark datasets and deduplicated by commit SHA to ensure each report maps to a unique ground-truth file set. Reports with missing ground-truth annotations or malformed commit references are discarded, resulting in a clean corpus of project-level bug metadata stored in a standardized Parquet format. In the second stage, source code embedding, all source files from each project repository are chunked into segments of 512 tokens with an overlap of 50 tokens, and each chunk is encoded using the \textbf{BAAI/bge-code-v1} \cite{bge_code} embedding model (768-dimensional, float32). These chunk-level embeddings are averaged to produce a single file-level embedding, stored as \texttt{float32} vectors in per-project pickle files. Bug report embeddings are generated analogously by encoding the full bug report text. The resulting per-project \textit{blob embedding databases} and \textit{bug metadata databases} serve as the foundation for both retrieval and domain gap computation.

\subsection{Evaluation Pipeline.}
The study adopts a \textbf{Retriever-Reranker} architecture. The retriever is a \textbf{FAISS IndexFlatIP} (inner-product search after L2 normalization), which retrieves the top-$K$ = 300 candidate source files for each bug report. The reranker is \textbf{TRANP-CNN} \cite{huo2019deep}, a CNN-based model that takes the (bug report, candidate file) pair as input and produces a relevance score used to re-rank the 300 candidates. This architecture decouples retrieval quality from reranking quality, enabling precise diagnosis of where failures occur.

\subsection{Training Scenarios.}
Every source$\rightarrow$target pair is evaluated under four training regimes, as described in Table~\ref{tab:scenarios}.

\begin{table*}[h]
\centering
\caption{Four evaluation scenarios and their data configurations.}
\label{tab:scenarios}
\begin{tabular}{|l|l|p{6.5cm}|}
\hline
\textbf{Scenario} & \textbf{Training Data} & \textbf{Research Relevance} \\ \hline
WP-small & 20\% target bugs only & Baseline: limited within-project training \\ \hline
WP-large & 80\% target bugs only & Upper bound: maximum within-project training \\ \hline
CP-cold-start & 100\% source, 0\% target & Zero-shot cross-project transfer \\ \hline
CP-transfer & 100\% source + 20\% target & Core CPL hypothesis: cross-project pre-training with minimal target fine-tuning \\ \hline
\end{tabular}
\end{table*}

\subsection{Evaluation Metrics}
Five complementary metrics are reported throughout this study to provide a complete picture of both retrieval and ranking performance. \textbf{Top-$K$ Recall} ($K \in \{1, 5, 10\}$) measures the fraction of bug reports for which the ground-truth file appears within the top $K$ ranked candidates. \textbf{Mean Average Precision (MAP)} measures the average precision at each rank position where a relevant file appears, rewarding models that place correct files higher in the ranking. \textbf{Mean Reciprocal Rank (MRR)} computes the mean of $1/\text{rank}$ across all bug reports, where rank is the position of the first correct answer. Together, Top-$K$ captures retrieval breadth while MAP and MRR capture ranking precision. It is important to note that these two families of metrics can diverge: a model may achieve high Top-10 recall but low MRR if the correct file is consistently ranked at position 9 rather than 1. Both perspectives are reported to avoid misleading conclusions.

\subsection{Domain Gap Measurement.}
To quantify the structural dissimilarity between project pairs, a \textbf{domain gap} score is computed for each source$\rightarrow$target combination by training a \texttt{LogisticRegression} binary classifier on the blob embeddings of the two projects. The rationale for this choice is three-fold. First, the \textbf{blob embeddings} (produced by BAAI/bge-code-v1) are high-dimensional, semantically rich representations that encode both vocabulary patterns and structural properties of source code files; a linear classifier applied to these embeddings directly tests whether the two projects are linearly separable in the shared embedding space, a necessary condition for domain shift to influence model transferability. Second, \textbf{logistic regression} is a well-understood, low-capacity classifier whose accuracy under balanced sampling is a clean measure of separability rather than a model-specific artifact: if the two projects are truly similar in embedding space, a logistic regression will achieve accuracy near 0.5 (chance); if they occupy disjoint regions, it will achieve accuracy near 1.0. Third, logistic regression has been widely used as a separability probe in domain adaptation literature and avoids the risk of a more powerful classifier overfitting on small embedding datasets. The resulting classification accuracy serves as the domain gap metric: a score of 1.0 indicates fully disjoint embedding distributions (high domain gap), while 0.5 indicates complete overlap (zero gap).

\section{Results and Discussion}

\subsection{Overall Scenario Performance (RQ1)}
Table~\ref{tab:overall} presents the mean performance across all 87 Python source$\rightarrow$target pairs for each evaluation scenario. A consistent ordering is observed across all five metrics:\textbf{ WP-large $>$ CP-transfer $>$ WP-small $>$ CP-cold-start}. CP-transfer, which uses the same 20\% target data budget as WP-small but augments it with 100\% source project pre-training, consistently outperforms WP-small while approaching WP-large. It is worth noting that existing SOTA CPBL results in the literature (e.g., TRANP-CNN \cite{huo2019deep}, FLIM \cite{liang2022modeling}, COOBA \cite{zhu2020cooba}) report MAP/MRR in the range of \textbf{0.3--0.5 for cross-project settings}, a range our best-performing pairs reach or exceed (see Showcase Pair B: MAP=0.770 for WP-large, MAP=0.703 for CP-transfer). The aggregate means are lower because they include hard large-codebase targets ($>$200K LoC) that depress the corpus-wide average.

\begin{table}[h]
\centering
\caption{Mean performance across all 87 Python pairs for each scenario.}
\label{tab:overall}
\begin{tabular}{|l|c|c|c|c|c|}
\hline
\textbf{Scenario} & \textbf{Top-1} & \textbf{Top-5} & \textbf{Top-10} & \textbf{MAP} & \textbf{MRR} \\ \hline
WP-large      & 0.075 & 0.191 & 0.252 & 0.220 & 0.246 \\ \hline
CP-transfer   & 0.065 & 0.183 & 0.244 & 0.199 & 0.227 \\ \hline
WP-small      & 0.056 & 0.161 & 0.215 & 0.177 & 0.202 \\ \hline
CP-cold-start & 0.004 & 0.021 & 0.040 & 0.039 & 0.042 \\ \hline
\end{tabular}
\end{table}

To assess whether the advantage of CP-transfer over WP-small is statistically reliable, a \textbf{Wilcoxon signed-rank test} is applied to each metric. This non-parametric paired test is appropriate because: (i) the observations are paired each source$\rightarrow$target pair contributes one CP-transfer and one WP-small result, and (ii) the distributions of improvements cannot be assumed normal across such a diverse project corpus. The one-sided alternative is tested (CP-transfer $>$ WP-small). Results are reported in Table~\ref{tab:wilcoxon}.

\begin{table}[h]
\centering
\caption{Wilcoxon signed-rank test: CP-transfer vs. WP-small, 87 Python pairs (one-sided, $H_1$: CP-transfer $>$ WP-small).}
\label{tab:wilcoxon}
\begin{tabular}{|l|c|c|c|}
\hline
\textbf{Metric} & \textbf{Win Rate} & \textbf{Mean Gain} & \textbf{$p$-value} \\ \hline
Top-1  & 36.5\% & +0.008 & 0.112 (n.s.) \\ \hline
Top-5  & 43.5\% & +0.020 & 0.019 ($\star$) \\ \hline
Top-10 & 47.1\% & +0.027 & 0.001 ($\star\star$) \\ \hline
MAP    & 63.5\% & +0.020 & 0.011 ($\star\star$) \\ \hline
MRR    & 67.1\% & +0.023 & 0.008 ($\star\star$) \\ \hline
\end{tabular}
\end{table}

The results demonstrate that CP-transfer produces statistically significant improvements over WP-small for MAP, MRR, Top-10, and Top-5 ($p \leq 0.02$), while Top-1 improvement is not statistically significant ($p = 0.112$). This finding reveals an important nuance: CPL's advantage is primarily in \textit{ranking quality} the model learns to order candidate files more accurately from cross-project training but does not consistently elevate the correct file to the very top position. Comparing CP-transfer against WP-large, the gap narrows considerably on recall metrics: CP-transfer matches or exceeds WP-large on Top-10 in over 50\% of pairs. The gap is more pronounced for MAP and MRR, indicating that WP-large retains a ranking precision advantage that CPL does not fully close with the current CNN-based architecture.

\subsection{Target Properties as Performance Predictors (RQ2)}
To identify the features most predictive of CPL performance, \textbf{Spearman rank correlation} is computed between project-level metadata features and CP-transfer performance for each metric. Spearman's $\rho$ is used rather than Pearson's because the project corpus spans several orders of magnitude in codebase size and bug count, making a linear association assumption inappropriate. A Spearman correlation of $\rho = +1.0$ indicates a perfect positive monotonic relationship, $\rho = -1.0$ a perfect negative monotonic relationship, and $\rho = 0$ indicates the absence of a systematic association. Results are reported in Table~\ref{tab:spearman}.

\begin{table*}[h]
\centering
\caption{Spearman rank correlations ($\rho$) between project features and CP-transfer performance (87 Python pairs). Significance: $\star\star\star$ ($p < 0.001$), $\star\star$ ($p < 0.01$), $\star$ ($p < 0.05$), n.s. (not significant).}
\label{tab:spearman}
\begin{tabular}{|l|c|c|c|c|}
\hline
\textbf{Feature} & \textbf{Top-10 $\rho$} & \textbf{MAP $\rho$} & \textbf{MRR $\rho$} & \textbf{Sig.} \\ \hline
Target LoC (tgt\_LoC)               & $-0.856$ & $-0.694$ & $-0.699$ & $\star\star\star$ \\ \hline
Target bug verbosity                & $+0.425$ & $+0.408$ & $+0.343$ & $\star\star\star$ \\ \hline
Source LoC (src\_LoC)               & $+0.407$ & $+0.311$ & $+0.321$ & $\star\star\star / \star\star$ \\ \hline
Target bug count (tgt\_n\_bugs)     & $+0.257$ & $+0.354$ & $+0.386$ & $\star / \star\star\star$ \\ \hline
Source bug count (src\_n\_bugs)     & $-0.006$ & $+0.024$ & $+0.064$ & n.s. \\ \hline
Domain gap                          & $< |0.12|$ & $< |0.12|$ & $< |0.12|$ & n.s. \\ \hline
\end{tabular}
\end{table*}

The most striking finding is that \textbf{target codebase size (tgt\_LoC) is by far the strongest predictor}, with $\rho = -0.856$ for Top-10 recall,a very strong negative correlation consistent across all metrics. As the target project's codebase grows, the correct source file is increasingly buried within a large candidate pool, making precise ranking progressively harder regardless of the source project used. The second-most important feature is target bug report verbosity ($\rho = +0.425$ for Top-10), which captures the descriptive richness of bug descriptions; more detailed reports provide stronger textual matching signals. Notably, source LoC also shows a significant positive correlation ($\rho = +0.407$ for Top-10): \textit{larger source projects carry richer and more diverse training signal}, enabling better generalisation. Critically, \textbf{source bug count and domain gap are both non-significant}: source bug count ($\rho = 0.064$ for MRR, n.s.) and domain gap ($|\rho| < 0.12$ for all metrics, n.s.) show no reliable predictive power. This finding directly challenges the intuitive assumption that source-target similarity is the primary driver of CPL success, and shifts the focus to target-project viability and report quality as the dominant factors.

Furthermore, the CPL gain (defined as CP-transfer$-$WP-small MRR) is \textbf{larger for data-scarce targets}: targets with fewer bug reports benefit proportionally more from cross-project pre-training, positioning CPL as most valuable precisely where it is most needed, for new, small, or low-activity projects with insufficient within-project training data.

\subsection{Source Project Selection (RQ3)}
To formally evaluate the feasibility of data-driven source project selection, \textbf{Kendall's rank correlation ($\tau$)} is computed between the ordering produced by candidate selection strategies and the oracle ordering (sorted by actual achieved CP-transfer MRR) for each of the 12 target projects with at least three available source candidates. Kendall's $\tau$ ranges from $-1$ (perfectly reversed ordering) to $+1$ (perfectly aligned ordering), with $\tau = 0$ indicating no systematic relationship between the heuristic and oracle rankings. Additionally, \textbf{Hit@1 accuracy} is reported: the fraction of targets for which the strategy selects the oracle-best source.

The \textbf{most-bugs heuristic}, selecting the source project with the largest bug corpus, achieves a Hit@1 of \textbf{41.7\%} (5 out of 12 targets) and a mean Kendall $\tau$ of \textbf{$+0.253$}. Table~\ref{tab:source_selection} compares strategies against the oracle upper bound and a random baseline.

\begin{table*}[h]
\centering
\caption{Source selection strategy comparison across 12 targets ($\geq$ 3 source candidates).}
\label{tab:source_selection}
\begin{tabular}{|l|c|c|c|c|}
\hline
\textbf{Strategy} & \textbf{Mean MRR} & \textbf{Mean MAP} & \textbf{Hit@1} & \textbf{Mean $\tau$ (MRR)} \\ \hline
Oracle (upper bound) & 0.320 & 0.273 & 100\% & --- \\ \hline
Most-bugs heuristic  & 0.293 & 0.249 & 41.7\% & $+0.253$ \\ \hline
Random (lower bound) & 0.262 & 0.228 & $\sim$15\% & $\sim$0.00 \\ \hline
\end{tabular}
\end{table*}

The most-bugs heuristic captures approximately 53\% of the gap between random and oracle selection ($(0.293 - 0.262) / (0.320 - 0.262) = 0.53$). Importantly, the domain gap and bug report similarity, both intuitively appealing proxies for source-target relevance, show no significant correlation with achieved CP-transfer performance ($\rho < 0.12$, $p > 0.05$ for all metrics). These findings establish that while a useful heuristic exists, source project selection remains a \textbf{partially solved problem}: the heuristic fails in 58\% of cases, and a learned meta-model capable of generalizing to unseen targets represents a clear direction for future work.

\subsection{FAISS Retrieval Ceiling (RQ4).}
Across all 94 source$\rightarrow$target pairs and all four scenarios, the fraction of bug reports for which the ground-truth file is absent from the FAISS top-300 candidates, denoted \textbf{pct\_unreachable}, is \textbf{0.00\%}. This finding has a critical architectural implication: every failure in the end-to-end pipeline is a \textit{reranking failure}, not a retrieval failure. The BAAI/bge-code-v1 embeddings are sufficiently expressive to retrieve the correct file within 300 candidates in all cases. The performance bottleneck lies entirely in TRANP-CNN's ability to promote that file from its initial FAISS rank (which can be as high as rank 167 for large-codebase targets) to within the top-10. Increasing the retrieval pool beyond $K = 300$ is therefore unlikely to yield meaningful improvements; instead, architectural advances in the reranker, such as replacing the CNN with a transformer-based or graph neural network-based model, represent the highest-impact research direction.

\subsection{Showcase Analysis: Three Illustrative Pairs}
To move beyond aggregate statistics and illustrate the range of CPL behaviours observed, three representative source$\rightarrow$target pairs are selected for in-depth analysis.

\textbf{Pair B: numpy $\rightarrow$ jupyterlab (Best CPL Performer).} This pair illustrates the strongest CPL performance in the corpus. Table~\ref{tab:pair_b} reports all five metrics across the four scenarios for the primary direction (numpy as source, jupyterlab as target), along with the FAISS retrieval baseline.

\begin{table}[h]
\centering
\caption{Pair B (numpy $\rightarrow$ jupyterlab): all metrics per scenario. jupyterlab: 39K LoC, 197 bug reports.}
\label{tab:pair_b}
\begin{tabular}{|l|c|c|c|c|c|}
\hline
\textbf{Scenario} & \textbf{Top-1} & \textbf{Top-5} & \textbf{Top-10} & \textbf{MAP} & \textbf{MRR} \\ \hline
WP-large      & 0.325 & 0.500 & 0.500 & \textbf{0.770} & \textbf{0.793} \\ \hline
CP-transfer   & 0.275 & 0.500 & 0.500 & 0.703 & 0.720 \\ \hline
WP-small      & 0.125 & 0.475 & 0.500 & 0.488 & 0.495 \\ \hline
CP-cold-start & 0.075 & 0.200 & 0.300 & 0.263 & 0.289 \\ \hline
FAISS baseline & ---  & ---   & ---   & ---  & 0.022 \\ \hline
\end{tabular}
\end{table}

CP-transfer (MAP=0.703, MRR=0.720) closely approaches WP-large (MAP=0.770, MRR=0.793) while using the same 20\% target data budget as WP-small (MAP=0.488, MRR=0.495). Both CP-transfer and WP-large achieve Top-10 recall of 0.500, matching the theoretical maximum for this test split. The FAISS baseline MRR of 0.022 underscores the magnitude of the reranker's contribution: over a \textbf{33$\times$ improvement} in ranking quality over raw embedding retrieval. These results are competitive with the cross-project SOTA range of MAP/MRR $\approx 0.3$--$0.5$ reported in the literature \cite{huo2019deep,liang2022modeling}.

\begin{figure*}[h]
\centering
\includegraphics[width=\textwidth]{els-cas-templates/figs/all_metrics_comparison.png}
\caption{Pair B (numpy $\leftrightarrow$ jupyterlab): all five metrics across four scenarios for both transfer directions. Top row: MRR with FAISS baseline overlay. Middle row: MAP. Bottom row: Top-K recall curves.}
\label{fig:pair_b}
\end{figure*}

Figure~\ref{fig:pair_b} presents the full metric breakdown for both directions of this pair. The forward direction (numpy $\rightarrow$ jupyterlab, left column) shows strong and consistent improvements across all scenarios. The reverse direction (jupyterlab $\rightarrow$ numpy) shows substantially lower performance, consistent with the commutativity analysis discussed below.

\textbf{Pair A: matplotlib $\leftrightarrow$ jupyterlab (Commutativity Contrast).} This pair illustrates the most pronounced commutativity asymmetry in the corpus. The two projects share an identical domain gap score of 0.972 (high dissimilarity), yet the CPL outcome reverses entirely when the source and target roles are swapped. Table~\ref{tab:pair_a} reports both MAP and MRR for both directions.

\begin{table*}[h]
\centering
\caption{Pair A: asymmetric CPL performance — both MAP and MRR, both directions.}
\label{tab:pair_a}
\begin{tabular}{|l|cc|cc|cc|cc|}
\hline
 & \multicolumn{2}{c|}{\textbf{WP-small}} & \multicolumn{2}{c|}{\textbf{WP-large}} & \multicolumn{2}{c|}{\textbf{CP-cold-start}} & \multicolumn{2}{c|}{\textbf{CP-transfer}} \\
\textbf{Direction} & MAP & MRR & MAP & MRR & MAP & MRR & MAP & MRR \\ \hline
matplotlib $\rightarrow$ jupyterlab & 0.472 & 0.480 & 0.587 & 0.589 & 0.048 & 0.050 & \textbf{0.665} & \textbf{0.667} \\ \hline
jupyterlab $\rightarrow$ matplotlib & \textbf{0.373} & \textbf{0.373} & 0.294 & 0.294 & 0.320 & 0.320 & 0.236 & 0.236 \\ \hline
\end{tabular}
\end{table*}

In the forward direction, where jupyterlab (39K LoC) serves as the target, CP-transfer (MAP=0.665, MRR=0.667) is the \textit{strongest} scenario, surpassing even WP-large (MAP=0.587). In the reverse direction, where matplotlib (249K LoC) serves as the target, CP-transfer (MAP=0.236, MRR=0.236) is the \textit{weakest} scenario, falling below WP-small (MAP=0.373). Critically, CP-cold-start is also competitive in the reverse direction (MAP=0.320 $>$ WP-large=0.294), suggesting that even zero-shot transfer is partially competitive when the target is large and hard. This reversal is entirely attributable to target codebase size, consistent with the dominant Spearman correlation of $\rho = -0.856$ for tgt\_LoC. Figure~\ref{fig:pair_a_comm} visualizes the commutativity asymmetry across both MRR and MAP, confirming this is not an artefact of a single metric.

\begin{figure*}[h]
\centering
\includegraphics[width=\textwidth]{els-cas-templates/figs/commutativity_contrast.png}
\caption{Pair A (matplotlib $\leftrightarrow$ jupyterlab): commutativity contrast for MRR (left) and MAP (right). The same domain gap (0.972) and same project pair produce opposite CPL outcomes depending on which project is the target, driven entirely by target codebase size.}
\label{fig:pair_a_comm}
\end{figure*}

\textbf{Pair C: numpy $\rightarrow$ scipy (Hard Target Ceiling).} scipy (438K LoC, 100 bug reports) represents the extreme large-codebase scenario. Table~\ref{tab:pair_c} shows that all four scenarios yield Top-1, Top-5, and Top-10 recall of 0.00, though MAP and MRR are non-zero (peaking at MAP=0.030, MRR=0.030 for WP-large), confirming that the reranker does learn to promote the GT file but not into the top-10. Despite this failure, pct\_unreachable = 0\%: the ground-truth file is present in the FAISS top-300 for every bug report. Figure~\ref{fig:pair_c_disp} displays the rank displacement distributions across scenarios, confirming that the reranker does improve rankings relative to FAISS but cannot overcome the scale of the search space (438K LoC, candidate pool of 300 from tens of thousands of files). This pair illustrates the target viability threshold: when the target codebase exceeds $\approx$200K LoC, no source project can compensate for the architectural limitation of the current CNN-based reranker.

\begin{table}[h]
\centering
\caption{Pair C (numpy $\rightarrow$ scipy): all metrics per scenario. scipy: 438K LoC — all Top-K recall is 0; MAP/MRR show small but nonzero signal.}
\label{tab:pair_c}
\begin{tabular}{|l|c|c|c|c|c|}
\hline
\textbf{Scenario} & \textbf{Top-1} & \textbf{Top-5} & \textbf{Top-10} & \textbf{MAP} & \textbf{MRR} \\ \hline
WP-large      & 0.000 & 0.000 & 0.000 & 0.030 & 0.030 \\ \hline
CP-transfer   & 0.000 & 0.000 & 0.000 & 0.024 & 0.024 \\ \hline
WP-small      & 0.000 & 0.000 & 0.000 & 0.004 & 0.004 \\ \hline
CP-cold-start & 0.000 & 0.000 & 0.000 & 0.006 & 0.006 \\ \hline
\end{tabular}
\end{table}

\begin{figure*}[h]
\centering
\includegraphics[width=\textwidth]{els-cas-templates/figs/rank_displacement_boxplot.png}
\caption{Pair C (numpy $\rightarrow$ scipy): rank displacement distributions per scenario. Positive displacement indicates the model improved the GT file's rank relative to FAISS. The model does improve ranks substantially (WP-large median $>$ 0), but the improvement is insufficient to break into the top-10 given scipy's 438K LoC search space.}
\label{fig:pair_c_disp}
\end{figure*}

\paragraph{Consolidated TRANP-CNN Results for Showcase Pairs.}
Table~\ref{tab:consolidated} provides a single-view reference of all five metrics across all four scenarios for the three showcase pairs (five directions total), enabling direct comparison across pairs and scenarios.

\begin{table*}[h]
\centering
\caption{TRANP-CNN results for all three showcase pairs — all 5 metrics $\times$ 4 scenarios. Values rounded to 3 d.p.}
\label{tab:consolidated}
\begin{tabular}{|l|l|c|c|c|c|c|}
\hline
\textbf{Pair / Direction} & \textbf{Scenario} & \textbf{Top-1} & \textbf{Top-5} & \textbf{Top-10} & \textbf{MAP} & \textbf{MRR} \\ \hline
\multirow{4}{*}{\shortstack[l]{Pair B\\numpy$\rightarrow$jupyterlab}}
 & WP-large      & 0.325 & 0.500 & 0.500 & 0.770 & 0.793 \\ \cline{2-7}
 & CP-transfer   & 0.275 & 0.500 & 0.500 & 0.703 & 0.720 \\ \cline{2-7}
 & WP-small      & 0.125 & 0.475 & 0.500 & 0.488 & 0.495 \\ \cline{2-7}
 & CP-cold-start & 0.075 & 0.200 & 0.300 & 0.263 & 0.289 \\ \hline
\multirow{4}{*}{\shortstack[l]{Pair A\\matplotlib$\rightarrow$jupyterlab}}
 & WP-large      & 0.150 & 0.500 & 0.500 & 0.587 & 0.589 \\ \cline{2-7}
 & CP-transfer   & 0.250 & 0.475 & 0.500 & 0.665 & 0.667 \\ \cline{2-7}
 & WP-small      & 0.125 & 0.425 & 0.500 & 0.472 & 0.480 \\ \cline{2-7}
 & CP-cold-start & 0.000 & 0.050 & 0.075 & 0.048 & 0.050 \\ \hline
\multirow{4}{*}{\shortstack[l]{Pair A (rev.)\\jupyterlab$\rightarrow$matplotlib}}
 & WP-small      & 0.054 & 0.081 & 0.162 & 0.373 & 0.373 \\ \cline{2-7}
 & CP-cold-start & 0.027 & 0.108 & 0.135 & 0.320 & 0.320 \\ \cline{2-7}
 & WP-large      & 0.027 & 0.081 & 0.162 & 0.294 & 0.294 \\ \cline{2-7}
 & CP-transfer   & 0.027 & 0.081 & 0.108 & 0.236 & 0.236 \\ \hline
\multirow{4}{*}{\shortstack[l]{Pair C\\numpy$\rightarrow$scipy}}
 & WP-large      & 0.000 & 0.000 & 0.000 & 0.030 & 0.030 \\ \cline{2-7}
 & CP-transfer   & 0.000 & 0.000 & 0.000 & 0.024 & 0.024 \\ \cline{2-7}
 & CP-cold-start & 0.000 & 0.000 & 0.000 & 0.006 & 0.006 \\ \cline{2-7}
 & WP-small      & 0.000 & 0.000 & 0.000 & 0.004 & 0.004 \\ \hline
\end{tabular}
\end{table*}

\section{Limitations}

This study is subject to several limitations that should be considered when interpreting the findings and applying the resulting guidelines.

\textbf{Single model evaluated in Phase 1.} The Phase 1 analysis is conducted exclusively with TRANP-CNN \cite{huo2019deep}. While this provides a controlled and internally consistent evaluation, the conclusions, particularly regarding the magnitude of CPL gain and the ranking precision gap relative to WP-large, may not generalize uniformly to other SOTA architectures such as FLIM \cite{liang2022modeling}, BL-GAN \cite{zhu2022bl}, or COOBA \cite{zhu2020cooba}, which differ in their reliance on structural code features (e.g., AST-based representations) and training objectives.

\textbf{Programming language coverage.} The current experimental corpus is predominantly Python-based, with a subset of Java projects. The generalizability of the Spearman correlations, particularly the dominant effect of tgt\_LoC, to C++, JavaScript, or polyglot projects remains empirically unvalidated. It is plausible that language-specific factors such as compilation unit granularity or file-level cohesion interact with codebase size in ways not captured by the current analysis.

\textbf{SOTA comparison scope.} The aggregate mean MRR of 0.20--0.25 reported in this study is computed over all 87 pairs, including very large-codebase targets that suppress the mean. Existing SOTA methods in the CPL literature (TRANP-CNN \cite{huo2019deep}, FLIM \cite{liang2022modeling}, COOBA \cite{zhu2020cooba}) report MAP/MRR in the range of \textbf{0.3--0.5 for cross-project settings}. Our best pairs (e.g., numpy$\rightarrow$jupyterlab: MAP=0.770) match or exceed this range. Importantly, TRANP-CNN was not originally evaluated in a within-project setting for Python projects; the original paper \cite{huo2019deep} reports results for Java projects under within-project evaluation, making direct numeric comparison inappropriate. All comparisons in this study are made against matched within-project baselines (WP-small, WP-large) and the FAISS retrieval baseline.

\textbf{Source selection heuristic limitations.} The most-bugs heuristic achieves Hit@1 = 41.7\%, meaning it selects a suboptimal source in the majority of cases. While this represents a positive Kendall $\tau$ and is better than random, it falls well short of oracle performance. The residual 47\% oracle gap suggests that project-level metadata alone is insufficient for reliable source selection, and that learned selection strategies trained on historical pair outcomes are needed.

\textbf{Data preprocessing and embedding quality.} All experiments rely on BAAI/bge-code-v1 embeddings generated from file-level chunking (512-token chunks with 50-token overlap). For very large source files that exceed a single chunk, the averaging of chunk embeddings may lose fine-grained structural information. Additionally, the preprocessing pipeline discards bug reports with missing or ambiguous ground-truth annotations, which may introduce a selection bias toward well-documented bugs and inflate retrieval recall compared to a less curated setting.



\section{Conclusion}
The conclusion goes here.


% Numbered list
% Use the style of numbering in square brackets.
% If nothing is used, default style will be taken.
%\begin{enumerate}[a)]
%\item 
%\item 
%\item 
%\end{enumerate}  

% Unnumbered list
%\begin{itemize}
%\item 
%\item 
%\item 
%\end{itemize}  

% Description list
%\begin{description}
%\item[]
%\item[] 
%\item[] 
%\end{description}  

% \clearpage %%Remove this from your manuscript

% Figure
% \begin{figure}%[]
%   \centering
% %    \includegraphics{}
%     \caption{}\label{fig1}
% \end{figure}


% \begin{table}%[]
% \caption{}\label{tbl1}
% \begin{tabular*}{\tblwidth}{@{}LL@{}}
% \toprule
%   &  \\ % Table header row
% \midrule
%  & \\
%  & \\
%  & \\
%  & \\
% \bottomrule
% \end{tabular*}
% \end{table}

% Uncomment and use as the case may be
%\begin{theorem} 
%\end{theorem}

% Uncomment and use as the case may be
%\begin{lemma} 
%\end{lemma}

%% The Appendices part is started with the command \appendix;
%% appendix sections are then done as normal sections
%% \appendix

% \section{}\label{}

% % To print the credit authorship contribution details
% \printcredits

%% Loading bibliography style file
%\bibliographystyle{model1-num-names}
\bibliographystyle{cas-model2-names}

% Loading bibliography database
\bibliography{cas-refs}

% Biography
%\bio{}
% Here goes the biography details.
%\endbio

%\bio{pic1}
% Here goes the biography details.
%\endbio

\end{document}

